
\documentclass{article}
\usepackage{colortbl}
\usepackage{makecell}
\usepackage{multirow}
\usepackage{supertabular}

\begin{document}

\newcounter{utterance}

\centering \large Interaction Transcript for game `ifeval', experiment `ifeval', episode 429 with refineonly{-}Qwen3.5{-}9B.
\vspace{24pt}

{ \footnotesize  \setcounter{utterance}{1}
\setlength{\tabcolsep}{0pt}
\begin{supertabular}{c@{$\;$}|p{.15\linewidth}@{}p{.15\linewidth}p{.15\linewidth}p{.15\linewidth}p{.15\linewidth}p{.15\linewidth}}
    \# & \multicolumn{2}{c}{Player} && \multicolumn{2}{c}{Game Master} \\
    \hline

    \theutterance \stepcounter{utterance}  
    & & & \multicolumn{4}{p{0.6\linewidth}}{
        \cellcolor[rgb]{0.9,0.9,0.9}{
            \makecell[{{p{\linewidth}}}]{
                \texttt{\tiny{[P1$\langle$GM]}}
                \texttt{Is "jiang" a Chinese name? What are other names similar to "jiang"? Separate your two answers with ******} \\
            }
        }
    }
     \\ \\

    \theutterance \stepcounter{utterance}  
    & \multicolumn{4}{p{0.6\linewidth}}{
        \cellcolor[rgb]{0.561,0.737,0.561}{
            \makecell[{{p{\linewidth}}}]{
                \texttt{\tiny{[P1$\rangle$GM]}}
                \texttt{Yes, "Jiang" (江) is a very common Chinese surname. It is the 15th most common surname in China and is also widely used in other East Asian countries like Korea and Japan (where it is often romanized as Kang or Ko). The character literally means "river," and many families adopted this surname historically because they lived near major rivers or to honor a river deity.} \\
\\ 
\texttt{******} \\
\\ 
\texttt{Other Chinese names that sound similar to "Jiang" or share the same pronunciation include:} \\
\texttt{*   **Jiang** (江): The standard surname.} \\
\texttt{*   **Jiang** (蒋): A common surname pronounced exactly the same, though the character means "to select" or "to choose."} \\
\texttt{*   **Jiang** (姜): Another frequent surname with the same pronunciation, often associated with ginger.} \\
\texttt{*   **Jiang** (将): A less common surname with the same sound, meaning "general" or "about to."} \\
\\ 
\texttt{Note: While these share the same pronunciation in Mandarin, they are distinct characters with different meanings and sometimes different stroke counts.} \\
            }
        }
    }
    & & \\ \\

    \theutterance \stepcounter{utterance}  
    & & & \multicolumn{2}{p{0.3\linewidth}}{
        \cellcolor[rgb]{0.9,0.9,0.9}{
            \makecell[{{p{\linewidth}}}]{
                \texttt{\tiny{[GM$|$GM]}}
                {'combination:two_responses': {}}
            }
        }
    }
    & & \\ \\

    \theutterance \stepcounter{utterance}  
    & & & \multicolumn{2}{p{0.3\linewidth}}{
        \cellcolor[rgb]{0.9,0.9,0.9}{
            \makecell[{{p{\linewidth}}}]{
                \texttt{\tiny{[GM$|$GM]}}
                \texttt{game\_result = WIN} \\
            }
        }
    }
    & & \\ \\

\end{supertabular}
}

\end{document}
